% !TEX root = main.tex
% !TEX TS-program = pdflatex
% !BIB program = biber


\lstset{language=C}

%Hier den Artikel einbeziehen, Verhalten der Softwarearchitekten deckt sich mit dem hier erforschten Nutzerverhalten bezüglich UML siehe Ablage in Quellen für Masterarbeit und Link: https://www.sciencedirect.com/science/article/pii/S0950584920300252


%Artikel beachten: BPMN nutzung findet vor allem in bestimmten Notationen statt, viele Bereiche der Sprache werden von Anwendern nicht oder sehr selten benötigt....unterstützt Punkte der Architekten bzgl Notation zu schwer zu erlernern bewzug auf diesen Artikel abgespeichert in Ablage, LINK https://www.emerald.com/bpmj/article/16/1/181/273804/Opportunities-and-constraints-the-current-struggle

\chapter{Diskussion der Ergebnisse}

Um die eingangs formulierten Forschungsfragen zu beantworten, erwies sich der qualitative Forschungsansatz (Kap. 3) als geeignetes Mittel, um das Erfahrungswissen der Softwarearchitekten zu ergründen. Die Anforderungen, die Softwarearchitekten an Modellierungssprachen haben, und welche Herausforderungen und Verbesserungspotentiale sie für bestehende Modellierungssprachen sehen, kann durch die qualitative Methodik bestmöglich dargestellt werden. Die Ergebnisse des vorherigen Kapitels lassen die folgenden Schlussfolgerungen zu.

Im Hinblick auf die erste Forschungsfrage "Welche Anforderungen haben Softwarearchitekten an Modellierungssprachen?"  werden im nächsten Abschnitt die zentralen Erkenntnisse dargestellt.



\section{Anforderungen an Modellierungssprachen}

Für die befragten Softwarearchitekten ist eine Modellierungssprache ein Werkzeug, das den Arbeitsalltag erleichtern soll. Sowohl die von den Architekten genannten Eigenschaften von Modellierungssprachen als auch die Aussagen zu tätigkeitsspezifischen Anforderungen, lassen Schlüsse auf die für die Softwarearchitekten wichtigste übergeordnete Anforderung an eine Modellierungssprache zu. Diese Anforderungen können folgendermaßen zusammengefasst werden:

Die Erleichterung der Kommunikation mit allen Anspruchsgruppen steht klar im Mittelpunkt der formulierten Anforderungen. 

Die Softwarearchitekten setzten mit ihren Antworten unterschiedliche Schwerpunkte, welcher Anforderungstyp für sie der jeweilig wichtigste ist. Das übergeordnete Ziel liegt aber in der Erleichterung der Kommunikation, in jedem der insgesamt zehn geführten Interviews ist die Betonung der Kommunikationsarbeit ein zentraler Aspekt. Modellierungssprachen werden als Werkzeug begriffen dieses Ziel zu erreichen.

Dass die Modellierungssprachen als Werkzeug begriffen werden, das von den Architekten zur Erleichterung ihrer Arbeit verwendet wird, zeigt sich auch in den Antworten zu den verschiedenen Anforderungstypen. Die meisten befragten Softwarearchitekten sehen die anwenderbezogenen Anforderungen als am wichtigsten an. Hier stehen die Einfachheit der Sprache, die gute Erlernbarkeit und eine gute Darstellbarkeit im Fokus. Diese von den Softwarearchitekten genannten Anforderungen spiegeln sich letztlich auch in den Wünschen und dem identifizierten Verbesserungspotential, auf das an späterer Stelle eingegangen wird, wider.

Die formalen Anforderungen werden von den Softwarearchitekten als zweitwichtigste Anforderungskategorie genannt. Die anwendungsbezogenen Anforderungen ist die am wenigsten wichtige Anforderungsart. Doch auch bei den Architekten, die die anwendungsbezogene Anforderungskategorie als wichtigste genannt haben, wurden diese Kategorien häufig als Beispiel dafür genutzt, um zu betonen, dass z. B. nur mit Einhalten auch der formalen Anforderungen die Kommunikation zwischen Architekten überhaupt möglich sei.\footnote {Vgl. elektronischer Anhang Interviewtranskript 3, Fundstelle 151, S. 6 }

Selbst der vollständige Verzicht auf Arbeit mit Modellierungssprachen wird damit begründet, dass die Modellierungssprachen nicht verwendet werden, weil diese die Kommunikation nicht erleichtern.
Betrachtet man weiter die Ergebnisse der Kategorie „A Tätigkeit der Softwarearchitekten“, zeichnet sich eine weitere Verbindung ab. Auch hier steht die Kommunikation der Modelle mit allen Anspruchsgruppen, ob Management, Entwicklung oder Fachbereiche im Vordergrund. Dies lässt sich sowohl in der Unterkategorie „A1 Aufgabenfeld der Softwarearchitekten“ als auch in der Unterkategorie „A5 Herangehensweise der Modellierung“ erkennen. 

Die unterschiedlichen Aspekte des Aufgabenfeldes der Softwarearchitekten sind zumindest teilweise mit der Rolle im Unternehmen bzw. der Behörde erklärbar, die der jeweilige Architekt einnimmt. Die im Ergebnisteil genannten Aufgaben geben insgesamt zu erkennen, dass die Tätigkeit eines Softwarearchitekten nur mit starken Kommunikationsfähigkeiten zu bewältigen sind, da hier eine Schnittstellenfunktion bzw. eine Multiplikatorfunktion eingenommen wird. Die Koordinationsfähigkeiten und die Managementfähigkeiten können als Untergruppe der Kommunikationsfähigkeit aufgefasst werden.

\section{Herausforderungen und Verbesserungspotentiale}
Die Herausforderungen und Verbesserungspotentiale, die von den Interviewteilnehmern identifiziert worden sind, stehen in direktem Zusammenhang mit der Antwort auf die erste Forschungsfrage. 

Die formulierten Limitationen und Herausforderungen der Unterkategorien C1 und C2 greifen das Thema Kommunikation erneut auf. Es werden häufig fehlende Praxistauglichkeit genannt, unzureichende Werkzeuge und eine zu hohe Komplexität der Sprache. Daneben sehen die Befragten auch Limitationen in der eigenen Rolle. Hier wird besonders der hohe Schulungsaufwand beim Erlernen angeführt sowie die herausfordernde korrekte Benutzung der Modellierungssprache. Auch hier wird explizit erwähnt, dass Modellierungssprachen die Notwendigkeit von Kommunikation nicht ersetzen können.\footnote{Vgl. elektronischer Anhang Interviewtrankskript 7, Fundstelle 376, S. 6 }
Bei den Werkzeugen im Speziellen sehen die Architekten viele Herausforderungen, von der Komplexität der Werkzeuge über proprietäre Software zu fehlenden Darstellungsmethoden. \footnote{Vgl. Anhang Unterkategorie C2 Herausforderungen}.

Die Nutzung der unterschiedlichen Werkzeuge der Architekten zeigt, dass ein großer Teil der Tätigkeit, das Sichtbarmachen der Anforderungen betrifft. Die Tabelle 4.2 (S. 36) zeigt auf, dass die Softwarearchitekten Werkzeuge zur graphischen Darstellung am häufigsten verwenden. Die Nutzung der einfachen Werkzeuge könnte das Bedürfnis nach besseren und mächtigeren Werkzeugen erklären. Hierbei spielt die Forderung nach Integrationsfähigkeit in bestehende Strukturen eine große Rolle. So wird Draw.io von den Architekten auch deswegen verwendet, weil sich das Zeichentool in das Collaboration Tool von Atlassian\footnote{In den Interviews wird RV-Wiki gesprochen, dieses ist mit dem Collaboration Tool von Atlassian erstellt worden} integrieren lässt. Dieses wird vom Arbeitgeber bzw. Beschäftigungsgeber der Architekten sowohl für das Projektmanagement als auch zu Dokumentationszwecken vorgeschrieben.

Die Softwarearchitekten sehen gerade bei der Werkzeugnutzung enormes Verbesserungspotential. Die gestellten Anforderungen und die identifizierten Herausforderungen könnten aus Sicht der Befragten mithilfe von Werkzeugverbesserungen gelöst oder zumindest reduziert werden. 

Die Werkzeuge sollten standardisierte Schnittstellen bieten, Versionierungssysteme verwenden und automatisierte Konsistenzprüfungen zur Verfügung stellen. Insbesondere diese Forderungen lassen sich als Hinweis auf bestehende Medienbrüche und Aktualitätsprobleme interpretieren. Der Bedarf an automatisierten Prüfungen und synchronisierten Modellen spricht dafür, dass Modellierung in der Praxis häufig von mangelnder Pflege oder unklaren Verantwortlichkeiten behindert wird.

Ein weiterer Befund betrifft die hohe Relevanz von Usability und Accessibility. Die wiederholt geäußerte Erwartung einer whiteboard-ähnlichen Nutzbarkeit sowie die Forderung nach reduzierter Komplexität deuten darauf hin, dass bestehende Modellierungswerkzeuge von Anwendern häufig als zu schwergewichtig oder wenig intuitiv wahrgenommen werden. Dies kann als ein möglicher Erklärungsansatz für die in der Praxis häufig beobachtete geringe Akzeptanz formaler Modellierung interpretiert werden. Die Ergebnisse legen nahe, dass die Gebrauchstauglichkeit einen entscheidenden Einfluss auf die tatsächliche Nutzung von Modellierungswerkzeugen hat.

Die Nutzbarkeit soll verbessert werden, von barrierefreier Software zu einfacherer Handhabung der Tools bis hin zu Sprachsteuerung der Software. Auch bei der Nutzung mehrerer Anwender gleichzeitig sehen die Befragten Verbesserungspotential. 

Neben den bestehenden Herausforderungen sehen aber auch viele Architekten keinen grundsätzlichen Handlungsbedarf bei den von ihnen genutzten Sprachen, so ist es wenig überraschend, dass die Verbesserungsvorschläge, die die Semantik oder die graphische Notation betreffen nur eine sehr geringe Bedeutung haben.  

Zusammengefasst lässt sich feststellen, dass die Herausforderungen und die Verbesserungspotentiale sich auf die Verbesserung der Modellierungssoftware fokussieren. Auch Teile der allgemeinen Verbesserungsvorschläge können mithilfe der Modellierungssoftware bzw. mit Anpassung der Modellierungswerkzeuge umgesetzt werden. Insgesamt lassen sich die Ergebnisse dahingehend interpretieren, dass der praktische Nutzen von Modellierungswerkzeugen weniger durch den Funktionsumfang im engeren Sinne bestimmt wird, sondern vielmehr durch ihre Integrationsfähigkeit in bestehende Arbeitsprozesse, ihre Benutzerfreundlichkeit sowie ihre Unterstützung kollaborativer und langfristiger Nutzungsszenarien. Damit verschiebt sich der Fokus von der Modellierung als isolierter Expertentätigkeit hin zu einer niedrigschwelligen, teamorientierten und in den Entwicklungsalltag integrierten Praxis.


\section{Nutzungsverhalten in der Praxis}
Ursächlich für die Nennung von UML und BPMN als am häufig genutzte Modellierungssprachen ist im Fall der befragten folgende Sachlage. Sowohl im Ausbildungsberuf des Fachinformatikers als auch im Studium der Informatik werden Modellierungssprachen behandelt. Im Falle der Ausbildung wird UML explizit genutzt, da dies integrale Bestandteil des Lehrplanes ist. Im Falle eines Studiums wird UML als theoretisches Wissen in einschlägigen Lehrbüchern behandelt.
Darüber hinaus gibt es für die deutsche Verwaltung/deutsche Behörden eine Architekturrichtlinie, die BPMN und UML als Modellierungssprache vorgibt.
\footnote{ IT-Planungsrat des Bundes,  https://www.it-planungsrat.de/fileadmin/beschluesse/2025/Beschluss_2025_17_Föderale_IT-Architekturrichtlinie_Nationale_IT-Architekturrichtlinie.pdf, aufgerufen am 26.02.2026}
Da die Interviewten aktuell alle im Behördenumfeld arbeiten, besteht die Anforderung, die in der IT-Architekturrichtlinie genannten Modellierungssprachen zu nutzen. 

\section{Spannungfelder}
Die meisten identifizierten Spannungsfelder können bei der Herangehensweise an die Modellierung und der Definition von Modellierungsende beobachtet werden. Hier treffen verschiedene Ansichten aufeinander. 

Während die Softwarearchitekten den Beginn der Modellierung oftmals ähnlich definiert haben als ersten Schritt bei der Betrachtung eines Problems oder der Anforderungssammlung oder dem Bestehen von Anforderungen, finden sich bei der Definition des Modellierungsendes deutlich unterschiedlichere Ansichten. Ein Einflussfaktor hierfür scheint die Position in der Organisation. So ist auffällig, dass die leitenden Softwarearchitekten ihr Modellierungsende deutlich früher sehen, im Extremfall mit Fertigstellen des Modells, als Softwarearchitekten, die selbst noch Entwickeln oder näher an der Entwicklungsabteilung sind. 

Auch bei der Herangehensweise an die Modellierung gibt es deutliche Differenzen. Hier trifft der völlige Verzicht auf den Einsatz von Modellierungssprachen auf die Einstellung, dass Entwicklung mit Modellierung gerade im Behördenumfeld besonders wichtig ist. Die Erklärung dieser enormen Lücke zwischen den Meinungen kann in den unterschiedlichen Lebensläufen der Softwarearchitekten und in der Spezialisierung des einzelnen Interviewteilnehmers liegen. Die meisten befragten Softwarearchitekten kommen bereits mit einem Softwareentwicklungshintergrund in die Softwarearchitektur und sehen Modellierungssprachen als notwendiges Werkzeug an. Die Gegenposition nimmt der Architekt ein, der als Datenbankadministrator in die Softwarearchitektur gekommen ist.

Auch bei der Verwendung von Checklisten und Mustern oder Best-Practices bei der Modellierung sind Gegensätze zu erkennen Während teilweise Softwarearchitekten mit Führungsverantwortung aktuell noch kaum oder wenige Vorlagen oder Muster für die Modellierung verwenden, sind sie dieser Vorgehensweise eher zugewandt und würden diese für die Zukunft gerne in ihren Abteilungen oder Teams einführen. Demgegenüber steht die Ansicht, dass gute Modellierung nur mit ausreichender Berufserfahrung als Softwarearchitekt und – darüber hinaus- davor als Entwickler möglich ist. Diese Seite argumentiert, dass aufgrund der ständigen und schnellen Veränderungen der Informationstechnologien, eine Nutzung von Mustervorlagen oder Checklisten nicht sinnvoll, nahezu unmöglich, sei.

Das Antwortverhalten ist in diesem Punkt für die Befragten inkonsistent. Denn es wird sich auch dafür ausgesprochen, dass es nur mit guten Einarbeitungsplänen möglich sei , neue Mitarbeiter in die Thematik einzuarbeiten. Auch dass jeder Softwarearchitekt Fortbildungen als wichtigen Teil der Tätigkeit erachtet, bildet hier einen Widerspruch. Fortbildungen können selbst als Muster betrachtet werden. In der Regel wird dort Erfahrungswissen und Theorie nach einem wiederkehrenden Schema den Teilnehmern beigebracht. Ebenfalls im Widerspruch stehen einzelne Aussagen, dass man mit den Modellierungssprachen mit dem heutigen Stand jede potentielle Anforderung modellieren kann. Dies lässt zumindest ein Muster auf abstrakter Ebene denkbar erscheinen.

Auch die formulierten Herausforderungen, gerade mit Blick auf die Probleme bei der Kommunikation, die von den Architekten einheitlich als großes Problem identifiziert worden sind, lassen erkennen, dass hier sich wiederholende ähnlich Problemstellungen im Arbeitsalltag entstehen. Eine Analyse dieser Probleme würde im Idealfall ein Schema entstehen lassen, wie man diesen Problemen in zukünftigen Projekten mit Modellierungstätigkeit begegnen kann.

Ein weiterer Gegensatz stellt die Einstellung der Befragten zum grundsätzlichen Aufbau der Modellierungssprachen dar, so gibt es Äußerungen, dass die Modellierungssprachen keine Verbesserungen benötigen, vor allem in den Punkten Semantik und graphische Notation. Diesen Aussagen entgegen stehen die Äußerungen darüber, dass eine Reduzierung der Komplexität der Sprache wünschenswert ist. Auch die genannten Limitationen, dass die Modellierungssprachen zu komplexe Beschreibungssprachen sind und nicht alles durch diese abbildbar ist, stellt sich direkt gegen die Äußerungen, dass die genutzten Modellierungssprachen keine Verbesserung benötigen. 

Zusammenfassend lässt sich feststellen, dass sich in den Antworten der Interviewteilnehmer einige Gegensätze finden. Die jeweils aktuelle Tätigkeit beeinflusst die Sicht auf jegliche Aspekte der Softwarearchitektur, aber nicht jede Inkonsistenz lässt sich mit dieser Begründung erklären, da die Antworten einzelner Architekten selbst in sich Widersprüche aufweisen.

\section{Einordnung in den Forschungsstand}

Vergleicht man die Erkenntnisse, die aus den Interviews gewonnen werden konnten mit den Forschungsergebnissen von Fettke, so lässt sich feststellen, dass dort ähnliche Beobachtungen aufgetreten sind. So hat Fettke bereits 2009 bei einer Untersuchung von 440 Anwendern festgestellt, dass die Anwender bei der Nutzung von Werkzeugen überwiegend Werkzeuge nutzen, die den Schwerpunkt „Zeichnen von Modellen“ haben. Weniger häufig werden andere Werkzeuge genutzt, die Prozessmanagement oder Softwareentwicklung als Hauptfokus haben. \footnote{Vgl. P. Fettke, Ansätze der Informationsmodellierung und ihre betriebswirtschaftliche Bedeutung: Eine Untersuchung der Modellierungspraxis in Deutschland. Schmalenbachs Z betriebswirtsch Forsch 61, 2009 S.565f  https://doi.org/10.1007/BF03373665} 

In dieser Studie wurden auch Problemfelder von Modellierungssprachen und Werkzeugen identifiziert. Vergleicht man die dort aufgeführten Erkenntnisse mit den Erkenntnissen aus dieser Forschungsarbeit, erkennt man auch hier diverse Übereinstimmungen. So wird auch von den dortigen Studienteilnehmern die Komplexität von Modellierungssprachen und eine unzureichende Ausdruckskraft genannt, mit dem Erfolgsfaktor leicht erlernbare Sprache. Bei den Werkzeugen wird die Komplexität, der Preis und das Ausfallrisiko des Werkzeuganbieters genannt. Es zeigt sich, die Werkzeuge werden von den dort Befragten anders betrachtet. Während in dieser Forschungsarbeit die Kostenfrage eines Werkzeuges nicht erscheint, spielt sie in Fettkes Studie eine große Rolle. 

Weitere aufgeführte Erfolgsfaktoren von Informationsmodellierung sind die gute Verständlichkeit der Modelle für Nicht-Informatiker und ausreichendes Training für die Modellierer. Auch in diesen Punkten gibt es Übereinstimmungen mit den Erkenntnissen aus dieser Forschungsarbeit.

Betrachtet man weitere wissenschaftliche Arbeiten, finden sich auch Überschneidungen mit den dortigen Erkenntnissen. So wird die Komplexität der Modellierungssprache BPMN in der Veröffentlichung von Recker ebenfalls bemängelt. Dieser Identifiziert 17 von den 50 nutzbaren Konstrukten als „Overhead“, die in der Praxis nicht oder sehr selten verwendet werden.
Weiter beschreibt die Veröffentlichung von Recker das Nutzungsverhalten der Experten folgendermaßen:

„Sie möchten ihre Arbeitsabläufe mit einfachen, grafischen Begriffen beschreiben.“
\footnote {Vgl. J. Recker, Opportunities and constraints: the current struggle with BPMN. Business Process Management Journal, Vol. 16 No. 1, 2010, S. 181–201, doi: https://doi.org/10.1108/14637151011018001 }

Dies lässt sich auch bei der Wahl der im Jahr 2010 angegebenen benutzten Werkzeuge erkennen. Hier führt Visio mit weitem Abstand an Nennungen. Auch diese Erkenntnis kann im Rahmen dieser Qualifikationsarbeit bestätigt werden, auch wenn bei den Interviewteilnehmern andere Werkzeuge bevorzugt werden. Die Erkenntnis hierbei scheint, dass das Nutzerverhalten sich auch über einen längeren Zeitraum nicht wesentlich verändert hat.

In der jüngsten Forschung wird bereits versucht die bekannten Probleme, die beim Modellieren bestehen, zu lösen. So kann man in der aktuellsten Veröffentlichung „The Practice of Modelling“ anlässlich der 18th IFIP Working Conference erkennen, dass ein Schwerpunkt die Verbesserung der Werkzeughandhabung ist. Schwantler et al. präsentieren hier ihre Ansätze zur Sprachsteuerung von Modellierungssoftware, um die Modellierung, hier mit UML, für Anwender zu vereinfachen. \footnote {Vgl. S. Schwantler, et al.,  Talk to Me! Toward Speech-Based UML Modeling. In: Fill, HG., Wautelet, Y., Ralyté, J., Zdravkovic, J. (eds) The Practice of Enterprise Modeling. PoEM 2025. Lecture Notes in Business Information Processing, vol 570. Springer, Cham., 2026, S. 83-101 https://doi.org/10.1007/978-3-032-12063-2_6}

Dieser Forschungsansatz trifft den Kern der Verbesserungsvorschläge dieser Forschungsarbeit, die Forderung nach einer sprachgesteuerten Software wurde bereits explizit erwähnt \footnote{Vgl. Anhang Interviewtranskript Softwarearchitekt 10, S. 4}.
Auch in dem Forschungsbeitrag von Kumarasinghe und Kirikova findet sich der von den Softwarearchitekten formulierte Wunsch nach Einbindung von KI in den Modellierungswerkzeugen wieder. Die Autoren stellen in ihrer Forschung einen Ansatz der Modellierungsunterstützung von Large Language Models bei BPMN Modellierung vor.\footnote{Vgl. A. Kumarasinghe, M. Kirikova, BPMN-Based Business Process Collaboration Modeling Using Large Language Models. In: Fill, HG., Wautelet, Y., Ralyté, J., Zdravkovic, J. (eds) The Practice of Enterprise Modeling. PoEM 2025. Lecture Notes in Business Information Processing, vol 570. Springer, Cham., 2026, S. 130-138  https://doi.org/10.1007/978-3-032-12063-2_9}

Besonders bemerkenswert ist die Übereinstimmung zwischen den Anforderungen der Anwender von vor 16 Jahren und denen aktueller Nutzer. Dies zeigt offenkundig, dass zentrale Problemstellungen bei Werkzeug-Usability trotz technologischer und methodischer Weiterentwicklungen bislang nicht gelöst wurden. 
Aktuelle Forschungsarbeiten zeigen jedoch, dass dieses Defizit inzwischen bereits erkannt wurde und auch als Verbesserungspotential adressiert wird. Zahlreiche aktuelle Veröffentlichungen befassen sich mit der Entwicklung von vielversprechenden Lösungsansätzen, die darauf abzielen, die identifizierten Herausforderungen zu überwinden.

Vergleicht man die Nutzung der Modellierungssprachen der Interviewteilnehmer mit der Untersuchung von Forsberg et al. (2025)\footnote{Vgl. A. Forsberg et al., Mind the Gaps: Which EA Modeling Notations Does Industry Actually Use? A Quantitative Investigation. In: Fill, HG., Wautelet, Y., Ralyté, J., Zdravkovic, J. (eds) The Practice of Enterprise Modeling. PoEM 2025. Lecture Notes in Business Information Processing, vol 570. Springer, 2026, S. 196ff, https://doi.org/10.1007/978-3-032-12063-2_13 } so erkennt man deutlich, dass Archimate und UML diejenigen Modellierungssprachen sind, die mit gleich vielen Nutzern den ersten Platz belegen. Danach wird BPMN als zweithäufigste genannt. Die Entity-Relationship-Diamme folgen wiederum danach.  Auch in dieser Studie wird festgestellt, dass die Anwender mehr als eine Modellierungssprache nutzen. Dies liegt, laut Forsberg, daran, dass die unterschiedlichen Arbeitgeber der Befragten, die sich aus dem privaten Unternehmen und der öffentliche Sektor zusammensetzen, mehr als eine Notation in ihrer Organisation verwenden. Die Studie zeigt weiterhin, dass die Befragten, die im öffentlichen Sektor tätig sind UML am meisten verwenden.

Die Erkenntnisse der Studie von Forsberg bezüglich der UML-Nutzung lassen sich auch in dieser Forschungsarbeit bestätigen. Da hier die Teilnehmer im öffentlichen Sektor tätig sind, findet auch hier UML die größte Anwendung. Auffällig ist jedoch der Unterschied im Nutzungsverhalten von Archimate. Hier finden sich in der Studie von Forsberg deutlich mehr Anwender, die auch im öffentlichen Sektor beschäftigt sind. Eine mögliche Erklärung hierzu könnte die höhere Anzahl der Teilnehmer, insgesamt 40, und die breitere Aufteilung über dreizehn Sektoren der Industrie und den öffentlichen Sektor sein. In der Studie selbst wird argumentiert, dass Archimate aufgrund seines nicht-domänenspezifischen Aufbaus flexibler einsetzbar ist und deshalb eine weitere Verbreitung findet.\footnote{Vgl. A. Forsberg et al., Mind the Gaps: Which EA Modeling Notations Does Industry Actually Use? A Quantitative Investigation a. a. O S. 197ff}

Sowohl die Studie als auch die hier gewonnen Forschungsergebnisse beruhen auf einer zu geringen Stichprobe um eine absolut allgemeingültige Aussage zu treffen, welche Modellierungssprachen am häufigsten verwendet werden. Beide Forschungsarbeiten konnten jedoch unabhängig voneinander feststellen, dass UML, Archimate und BPMN zu den meist genutzten Sprachen gehören.

In der Forschung zu Modellierungssprachen haben Bjeković et al. (2014) die These aufgestellt, dass es sinnvoller ist den grundlegenden Sinn und die Intention der Modellierer und von Modellierung zu erfassen, bevor man mit immer kleinteiligeren Modellierungssprachen-Erweiterungen versucht die speziellen Anforderungen der einzelnen Anwendergruppen zu lösen. Der These nach wird  der Zweck von Modellierungssprachen durch zwei Funktionen von Modellierungssprachen abgebildet: die linguistische Funktion und die Repräsentationsfunktion. Erstere soll dabei unterstützen die Domäne und die Vorstellung des Anwenders zur Domäne zu formen. Die zweite Funktion soll dabei helfen die Domäne mit einem zweckorientierten Modell darzustellen.
\footnote {Vgl. M Bjeković et al., Embracing Pragmatics. In: Yu, E., Dobbie, G., Jarke, M., Purao, S. (eds) Conceptual Modeling. ER 2014. Lecture Notes in Computer Science, vol 8824.  Springer, 2014, S. 431-444 https://doi.org/10.1007/978-3-319-12206-9_37}

Dieser Argumentation kann insoweit gefolgt werden, als dass anhand der vorliegenden Forschungsarbeit und den Beschreibungen der befragten Softwarearchitekten zur Herangehensweise und ihren Anforderungen an Modellierungssprachen, sich die dort erhaltenen Ergebnisse in diese zwei von Bjeković et al. entwickelten Funktionen einordnen lassen. Inwiefern diese Einordnung tatsächlich die (Weiter-)Entwicklung von Modellierungssprachen zu verbessern vermag oder die tatsächlichen Anforderungen der Softwarearchitekten an die Modellierungssprachen erfüllen kann, ist allerdings in Frage zu stellen.

Die in dieser Forschungsarbeit ausgearbeiteten Ergebnisse lassen eher den Schluss zu, dass die bestehenden Herausforderungen mit einer pragmatischen Sicht auf die Anwenderbedürfnisse, also der Verbesserung der Modellierungswerkzeuge, gelöst werden können. 

Da mit Veränderung der technischen Möglichkeiten auch eine Veränderung der Nutzeranforderungen einhergeht, ist der Beitrag dieser Forschungsarbeit jedoch auch, die aktuellen Anforderungen und Verbesserungsvorschläge der Anwender festzuhalten.

Anhand der aktuellsten Forschungsveröffentlichungen lässt sich erkennen, dass Modellierungssprachen und ihre praxisorientiere Anwendung weiter im Fokus der Forschung stehen. Die Entwicklungen im Bereich künstlicher Intelligenz haben hier eine neue Perspektive eröffnet, um die bestehenden Probleme zu bewältigen. Viele dieser Herausforderungen sind bereits in älteren Forschungsarbeiten erkannt worden. 



\section{Limitationen und kritische Reflexion}
Die vorliegende Untersuchung unterliegt mehreren Einschränkungen, die bei der Interpretation und Einordnung der Ergebnisse berücksichtigt werden müssen.
Die Aussagekraft der Ergebnisse ist durch die Größe und Zusammensetzung der Stichprobe naturgemäß begrenzt. Die Anzahl der befragten Experten ist mit zehn Interviewteilnehmern vergleichsweise klein, was bei qualitativen Studien zwar üblich ist, jedoch zugleich eine statistische Verallgemeinerbarkeit der hier gefundenen Erkenntnisse einschränkt. Darüber hinaus sind die Teilnehmer aktuell alle für die gleiche Behörde tätig und in diesem Arbeitsumfeld unterliegen sie gleichen Anforderungen, z.B. der Nutzung bestimmter Modellierungssprachen. Diese Homogenität kann dazu führen, dass die Ergebnisse nur eingeschränkt auf andere Kontexte übertragbar sind. Diese Einschränkung wird durch die unterschiedlichen Lebensläufe und Berufserfahrung zwar abgemildert, aber nicht vollständig ausgeräumt. 

Die qualitative Inhaltsanalyse ermöglicht zwar eine differenzierte und kontextnahe Auswertung der Interviewdaten, ist jedoch mit methodischen Grenzen verbunden. Die Interpretation der Aussagen eröffnet zwangsläufig Spielräume, wodurch subjektive Einflüsse des Forschers nicht vollständig ausgeschlossen werden können. Es besteht das Risiko, dass einzelne Aussagen unterschiedlich gewichtet oder interpretiert werden. Zudem erlaubt die Methode keine quantitative Bewertung der Relevanz einzelner Anforderungen im statistischen Sinne. Die Ergebnisse zeigen daher vor allem thematische Muster und inhaltliche Schwerpunkte, jedoch keine repräsentativen Häufigkeiten oder Prioritäten im engeren quantitativen Verständnis.

Auch im Rahmen der Datenerhebung ergeben sich potenzielle Limitationen. Die Dauer der Interviews sowie die Struktur des Leitfadens beeinflussen die erhobenen Informationen. Der zeitliche Rahmen der Interviews war mit etwa einer Stunde angesetzt, den Interviewteilnehmern wurde jedoch mitgeteilt, dass keine weiteren terminlichen Zwänge von Seiten des Forschers bestehen. Dennoch ist nicht vollständig auszuschließen, dass bestimmte Aspekte nicht ausführlich diskutiert wurden.

Es ist ebenfalls nicht auszuschließen, dass bei manchen Begrifflichkeiten ein unterschiedliches Verständnis vorliegt. Zum einen zwischen Forscher und Befragten, aber auch die Softwarearchitekten untereinander können ein unterschiedliches Begriffsverständnis haben. Gerade bei den Anforderungsbegrifflichkeiten, die selbst unter Forschern (vgl. Kapitel 2) nicht vollständig einheitlich definiert sind, besteht zwangsläufig die Möglichkeit von Interpretationsspielräumen im Interviewkontext.

Eine weitere Limitation besteht im Hinblick auf das Ziel der theoretischen Sättigung. Diese entsteht in der Theorie erst, wenn weiterer Forschungsinput, hier Interviews, keine neuen Erkenntnisse mehr hervorbringen.  Dies ist in dieser Forschungsarbeit nicht vollständig aufgetreten. Zwar sind in den Interviews viele gleiche oder ähnliche Aspekte aufgetreten, gerade im Hinblick auf Anforderungen im Kontext von Kommunikation oder Verbesserungsmöglichkeiten. Doch auch im letzten Interview sind neue Erkenntnisse zu Verbesserungsmöglichkeiten wie die bereits oben genannte Sprachsteuerung erwähnt worden. Dies lässt die Option offen, dass weitere Interviews noch weitere neue Erkenntnisse zum Vorschein gebracht hätten. Aufgrund des begrenzten zeitlichen Rahmens dieser Forschungsarbeit konnte daher die theoretische Sättigung nicht vollständig erreicht werden. 

Auch die Verwendung des Leitfadens kann die Gesprächsführung beeinflussen und unbeabsichtigt dem Interviewpartner thematische Schwerpunkte vorgeben, oder durch im Gesprächsfluss entstandene Suggestivfragen die Meinung des Befragten beeinflussen. Die gestellten Suggesitivfragen wurden aber in der qualitativen Analyse weitestgehend als nicht verwendbar erachtet und bei der Auswertung nicht berücksichtigt.

Darüber hinaus besteht bei Experteninterviews grundsätzlich die Möglichkeit sozial erwünschter Antworten, etwa wenn Teilnehmende ihre Arbeitsweise oder organisatorische Rahmenbedingungen in einem besonders positiven Licht darstellen oder sich an vermeintlichen Best Practices orientieren. Die Tatsache, dass der Forscher auch Arbeitskollege der befragten Experten ist, kann das Antwortverhalten ebenfalls beeinflusst haben. Dagegen spricht jedoch, dass es zu keinem der Interviewteilnehmer ein direktes hierarchisches Über- oder Unterordnungsverhältnis gibt. Die klare Kommunikation, dass die Interviews vertraulich und nur zu Forschungszwecken verwendet werden, ebenso die zugewandte Interviewführung, sollten eine freie und offene Gesprächsatmosphäre erzeugen. Die in den Interviews auftauchende offene Kritik an der eignen Tätigkeit oder andere sensible Antworten, sprechen gegen ein angepasstes Antwortverhalten. Das Risiko lässt sich allerdings nicht vollständig ausräumen.

Ein weiterer möglicher Einflussfaktor ergibt sich aus der Vorerfahrung der Experten mit bestimmten Modellierungssprachen. Diese Vertrautheit kann die Wahrnehmung und Formulierung von Anforderungen und Verbesserungsvorschlägen geprägt haben, indem bekannte Konzepte, Darstellungsformen oder Denkweisen als Referenzrahmen dienten. Infolgedessen könnten Anforderungen eher inkrementell aus bestehenden Ansätzen abgeleitet worden sein, anstatt vollständig unabhängig formuliert zu werden. Da fast alle Experten UML verwenden, teilweise seit der Ausbildung, ist nicht auszuschließen, dass diese Modellierungssprache als „gutes Beispiel“ verankert ist und die dort gelernten Modellierungskenntnisse mit anderen Modellierungssprachen abgeglichen wird.  Da es sich aber um Experteninterviews handelt, ist gerade die Vorerfahrung und die Tätigkeit der Grund, weshalb der Interviewpartner ausgewählt wurde. 

\section{Implikationen für Praxis und Forschung}
Die Ergebnisse dieser Arbeit liefern sowohl für die Praxis als auch für die weitere Forschung relevante Ansatzpunkte. Insbesondere im Hinblick auf die identifizierten Anforderungen an Modellierungswerkzeuge ergeben sich konkrete Hinweise für Anbieter von Modellierungssoftware. Die herausgearbeiteten Aspekte wie Schnittstellentauglichkeit, automatisierte Generierung von Modellen, Barrierefreiheit, KI-Assistenz und Kollaborationsfähigkeit können als Orientierung für die Weiterentwicklung bestehender Produkte dienen, indem Funktionalitäten gezielt an den Bedürfnissen der Anwender ausgerichtet und bestehende Defizite adressiert werden. Auf diese Weise kann die Gebrauchstauglichkeit sowie die Akzeptanz entsprechender Werkzeuge in der Praxis verbessert werden.

Neben den praktischen Implikationen ergeben sich auch Ansatzpunkte für weitere Forschung. Insbesondere erscheint eine vertiefende empirische Untersuchung der tatsächlichen Nutzung von Modellierungswerkzeugen in unterschiedlichen Anwendungskontexten sinnvoll. Ziel solcher Studien könnte es sein Nutzungsmuster, Herausforderungen sowie Erfolgsfaktoren im praktischen Einsatz genauer zu analysieren und die hier gewonnenen Ergebnisse weiter zu validieren und zu differenzieren. Zudem könnte eine Ausweitung der Stichprobe oder die Betrachtung weiterer Domänen dazu beitragen die Übertragbarkeit der Ergebnisse zu erhöhen. 

Darüber hinaus bilden die vorliegenden Erkenntnisse eine geeignete Grundlage für weiterführende Entwicklungsarbeiten. In zukünftigen Forschungsprojekten könnten Modellierungswerkzeuge konzipiert und prototypisch umgesetzt werden, welche die identifizierten Anforderungen systematisch berücksichtigen und deren Nutzen anschließend empirisch überprüft wird. 

Die vorangegangene Diskussion hat die Ergebnisse dieser Arbeit eingeordnet und kritisch reflektiert. Im folgenden Kapitel werden die wesentlichen Erkenntnisse der Forschungsarbeit zusammengefasst.
